\ifdefined\FULLARTICLE
\else
  \documentclass[12pt,a4paper]{article}
  \usepackage[T2A]{fontenc}
  \usepackage[utf8]{inputenc}
  \usepackage[russian]{babel}
  \usepackage{amsmath, amssymb, amsthm}
  \usepackage{geometry}
  \usepackage{booktabs} % для красивых таблиц
  \usepackage{tabularx}
  \geometry{margin=2.5cm}

  \theoremstyle{definition}
  \newtheorem{remark}{Remark}

  \begin{document}

  \begin{center}
    \Large\textbf{Examples of Productivity Calculations with AI}
  \end{center}
\fi

\ifdefined\FULLARTICLE
  \newcommand{\detailedexampleheading}[1]{\subsubsection{#1}}
  \newcommand{\detailedsummaryheading}[1]{\subsection{#1}}
\else
  \newcommand{\detailedexampleheading}[1]{\subsection*{#1}}
  \newcommand{\detailedsummaryheading}[1]{\section*{#1}}
\fi

\detailedsummaryheading{Elementary M0--M1 Examples}

Сначала рассмотрим четыре самостоятельных расчётных примера. Во всех случаях базовая трудоёмкость единицы работы принимается равной $X = 4$~часа.

Идеализированное время выполнения с ИИ вычисляется по формуле:
\[
  T_{ai}^{(0)} = \frac{1}{P} \sum_{i=1}^{N} Z_i k_i X.
\]

Условие ускорения в идеализированной модели можно записать через порог параллелизма:
\[
  P > \frac{\sum_{i=1}^{N} Z_i k_i}{\sum_{i=1}^{N} Z_i} = \bar{k}_w.
\]
Если доступный параллелизм $P$ не превышает порог $\bar{k}_w$, ускорение не достигается. Условие $P > \bar{k}_w$ необходимо во всех вариантах модели; достаточным оно является только в идеализированной модели идеально делимых задач.

Для сопоставления с более реалистичной моделью неделимых задач дополнительно используется нижняя оценка времени завершения всех работ:
\[
  T_{ai} \ge \max\left\{ \frac{1}{P} \sum_{i=1}^{N} Z_i k_i X,\; \max_{i \in \overline{1,N}} (Z_i k_i X) \right\}.
\]

Расчёты выполнены при $C(P) \equiv 1$ и $\gamma(P) \equiv 1$. Примеры 1--4 относятся к уровням M0, где работа считается идеально делимой, и M1, где учитываются независимые неделимые задачи. В иллюстративном сценарии дополнительно рассматриваются уровень M2 с графом зависимостей и уровень M4 с человеком как последовательным ресурсом. Для каждой задачи на уровне M4 выделяется человеческая составляющая $h_i$ в разложении $k_i = a_i + h_i$, вычисляются суммарное человеческое время $H = X \sum_i Z_i h_i$ и потолок ускорения $1/\bar{h}_w$. Делегированная задача разбивается на входную спецификацию, автономную агентную фазу и итоговую проверку. С момента постановки до приёмки она удерживает назначенного агента, в том числе при ожидании свободного разработчика. Расписания проверяются на отсутствие пересечений человеческих интервалов и переиспользования заблокированных потоков; чувствительность к $\gamma(P)>1$ обсуждается при сравнении вариантов. Уровень M3, описывающий межагентные координационные издержки $C(P)$, в расчётах не используется.

\detailedexampleheading{Example 1}

\paragraph{Parameters.}
\begin{table}[ht!]
\centering
\caption{Параметры примера 1}
\label{tab:example-1-params}
\begin{tabular}{lc}
\toprule
Параметр & Значение \\
\midrule
$N$ (число задач) & 3 \\
$\{Z_i\}$ (сложности) & $\{2, 3, 1\}$ \\
$\{k_i\}$ (коэфф.\ ИИ) & $\{0.5, 0.7, 0.4\}$ \\
$P$ (параллелизм) & 2 \\
$X$ (базовая трудоёмкость) & 4 ч \\
\bottomrule
\end{tabular}
\end{table}

\paragraph{Calculations.}
\begin{align*}
  \sum Z_i &= 2 + 3 + 1 = 6, \\
  \sum Z_i k_i &= 2{\cdot}0.5 + 3{\cdot}0.7 + 1{\cdot}0.4 = 3.5, \\
  T_h &= 6 \cdot 4 = 24~\text{ч}, \\
  T_{ai}^{(0)} &= \frac{1}{2} \cdot 3.5 \cdot 4 = 7~\text{ч}, \\
  S_E &= \frac{24}{7} \approx 3.43.
\end{align*}

Для модели неделимых задач веса задач с ИИ равны:
\[
  \{Z_i k_i X\} = \{4.0, 8.4, 1.6\}~\text{ч}.
\]
При оптимальном распределении первый исполнитель получает задачу длительностью $8.4$~ч, а второй --- две задачи общей длительностью $4.0 + 1.6 = 5.6$~ч. Поэтому достижимое время завершения равно $8.4$~ч, а ускорение составляет
\[
  S_E^{\text{ach}} = \frac{24}{8.4} \approx 2.86.
\]

\paragraph{Interpretation.}
Во всех трёх задачах $k_i<1$, поэтому ИИ снижает их трудоёмкость, а параллелизм $P=2$ усиливает этот эффект. Идеализированная модель даёт ускорение $\times 3.43$. После учёта неделимости задач оно снижается до $\times 2.86$, поскольку время завершения ограничено самой длинной задачей продолжительностью $8.4$~ч.

\detailedexampleheading{Example 2}

\paragraph{Parameters.}
\begin{table}[ht!]
\centering
\caption{Параметры примера 2}
\label{tab:example-2-params}
\begin{tabular}{lc}
\toprule
Параметр & Значение \\
\midrule
$N$ (число задач) & 4 \\
$\{Z_i\}$ (сложности) & $\{5, 2, 4, 3\}$ \\
$\{k_i\}$ (коэфф.\ ИИ) & $\{0.8, 0.6, 0.9, 0.5\}$ \\
$P$ (параллелизм) & 3 \\
$X$ (базовая трудоёмкость) & 4 ч \\
\bottomrule
\end{tabular}
\end{table}

\paragraph{Calculations.}
\begin{align*}
  \sum Z_i &= 5 + 2 + 4 + 3 = 14, \\
  \sum Z_i k_i &= 5{\cdot}0.8 + 2{\cdot}0.6 + 4{\cdot}0.9 + 3{\cdot}0.5 = 10.3, \\
  T_h &= 14 \cdot 4 = 56~\text{ч}, \\
  T_{ai}^{(0)} &= \frac{1}{3} \cdot 10.3 \cdot 4 \approx 13.73~\text{ч}, \\
  S_E &= \frac{56}{13.73} \approx 4.08.
\end{align*}

Для модели неделимых задач веса задач с ИИ равны:
\[
  \{Z_i k_i X\} = \{16.0, 4.8, 14.4, 6.0\}~\text{ч}.
\]
При оптимальном распределении нагрузки исполнители получают соответственно $16.0$, $14.4$ и $6.0 + 4.8 = 10.8$~ч работы. Достижимое время завершения равно $16.0$~ч, а ускорение составляет
\[
  S_E^{\text{ach}} = \frac{56}{16.0} = 3.50.
\]

\paragraph{Interpretation.}
В примере 2 для всех задач выполняется $k_i<1$, а коэффициент параллелизма равен $P=3$. Идеализированное ускорение $\times 4.08$ оказывается наибольшим среди первых трёх примеров. После учёта неделимости время завершения ограничивает задача длительностью $16.0$~ч, поэтому достижимое ускорение снижается до $\times 3.50$.

\detailedexampleheading{Example 3}

\paragraph{Parameters.}
\begin{table}[ht!]
\centering
\caption{Параметры примера 3}
\label{tab:example-3-params}
\begin{tabular}{lc}
\toprule
Параметр & Значение \\
\midrule
$N$ (число задач) & 3 \\
$\{Z_i\}$ (сложности) & $\{8, 5, 6\}$ \\
$\{k_i\}$ (коэфф.\ ИИ) & $\{1.1, 0.7, 0.8\}$ \\
$P$ (параллелизм) & 2 \\
$X$ (базовая трудоёмкость) & 4 ч \\
\bottomrule
\end{tabular}
\end{table}

\paragraph{Calculations.}
\begin{align*}
  \sum Z_i &= 8 + 5 + 6 = 19, \\
  \sum Z_i k_i &= 8{\cdot}1.1 + 5{\cdot}0.7 + 6{\cdot}0.8 = 17.1, \\
  T_h &= 19 \cdot 4 = 76~\text{ч}, \\
  T_{ai}^{(0)} &= \frac{1}{2} \cdot 17.1 \cdot 4 = 34.2~\text{ч}, \\
  S_E &= \frac{76}{34.2} \approx 2.22.
\end{align*}

Для модели неделимых задач веса задач с ИИ равны:
\[
  \{Z_i k_i X\} = \{35.2, 14.0, 19.2\}~\text{ч}.
\]
При оптимальном распределении один исполнитель получает задачу длительностью $35.2$~ч, а второй --- две задачи общей длительностью $19.2 + 14.0 = 33.2$~ч. Достижимое время завершения равно $35.2$~ч, а ускорение составляет
\[
  S_E^{\text{ach}} = \frac{76}{35.2} \approx 2.16.
\]

\paragraph{Interpretation.}
В примере 3 ускорение сохраняется, хотя для одной задачи ИИ увеличивает трудоёмкость ($k_i = 1.1$). Эта же задача длительностью $35.2$~ч ограничивает расписание. Поэтому идеализированное ускорение $\times 2.22$ лишь немного превышает достижимое значение $\times 2.16$; оба результата ниже, чем в примерах 1 и 2.

\detailedexampleheading{Example 4}

\paragraph{Parameters.}
\begin{table}[ht!]
\centering
\caption{Параметры примера 4}
\label{tab:example-4-params}
\begin{tabular}{lc}
\toprule
Параметр & Значение \\
\midrule
$N$ (число задач) & 3 \\
$\{Z_i\}$ (сложности) & $\{4, 3, 3\}$ \\
$\{k_i\}$ (коэфф.\ ИИ) & $\{2.4, 2.1, 2.3\}$ \\
$P$ (параллелизм) & 2 \\
$X$ (базовая трудоёмкость) & 4 ч \\
\bottomrule
\end{tabular}
\end{table}

\paragraph{Calculations.}
\begin{align*}
  \sum Z_i &= 4 + 3 + 3 = 10, \\
  \sum Z_i k_i &= 4{\cdot}2.4 + 3{\cdot}2.1 + 3{\cdot}2.3 = 22.8, \\
  \bar{k}_w &= \frac{22.8}{10} = 2.28, \\
  T_h &= 10 \cdot 4 = 40~\text{ч}, \\
  T_{ai}^{(0)} &= \frac{1}{2} \cdot 22.8 \cdot 4 = 45.6~\text{ч}, \\
  S_E &= \frac{40}{45.6} \approx 0.88.
\end{align*}

Пороговое условие ускорения не выполняется:
\[
  P = 2 < \bar{k}_w = 2.28.
\]
Следовательно, доступного параллелизма недостаточно, чтобы компенсировать рост трудоёмкости задач при использовании ИИ.

Для модели неделимых задач веса задач с ИИ равны:
\[
  \{Z_i k_i X\} = \{38.4, 25.2, 27.6\}~\text{ч}.
\]
При $P=2$ оптимальное распределение имеет вид $\{38.4\}$ и $\{25.2 + 27.6 = 52.8\}$~ч. Поэтому достижимое время завершения равно $52.8$~ч, а ускорение оказывается ниже идеализированной оценки:
\[
  S_E^{\text{ach}} = \frac{40}{52.8} \approx 0.76 < 1.
\]

Нижняя оценка времени завершения равна $\max\{45.6,\; 38.4\} = 45.6$~ч, но не достигается: оптимальное значение составляет $52.8$~ч. Это единственный из четырёх примеров, в котором выражение $\max\{W/P,\; \max_i Z_i k_i X\}$ строго меньше оптимума. Тем самым видно, что оно задаёт нижнюю границу, а не гарантированно достижимый результат. Расхождение остаётся в пределах гарантии Грэма: $52.8 \le (2 - 1/2) \cdot 45.6 = 68.4$~ч.

Достаточный критерий ускорения также не выполняется. Доля самой длинной задачи равна $\mu = M_N/T_h = 38.4/40 = 0.96$, а порог --- $P - (P-1)\,\mu = 2 - 0.96 = 1.04$, тогда как $\bar{k}_w = 2.28 > 1.04$. Более того, нарушено даже необходимое условие $\bar{k}_w < P = 2$, что согласуется с отсутствием ускорения.

\paragraph{Interpretation.}
В примере 4 все коэффициенты $k_i$ существенно больше единицы, поэтому ИИ заметно увеличивает трудоёмкость каждой задачи. При $\bar{k}_w = 2.28$ доступного параллелизма $P=2$ недостаточно, чтобы компенсировать это замедление. Даже идеализированная модель показывает рост времени с $40$ до $45.6$~ч, а после учёта неделимости и неравномерной нагрузки время завершения увеличивается до $52.8$~ч.

\detailedsummaryheading{Summary Comparison}

\begin{table}[ht!]
\centering
\caption{Сводные результаты расчёта}
\label{tab:summary-results}
\begin{tabular}{lcccc}
\toprule
Показатель & Пример 1 & Пример 2 & Пример 3 & Пример 4 \\
\midrule
$T_h$ (без ИИ), ч & 24.0 & 56.0 & 76.0 & 40.0 \\
$T_{ai}^{(0)}$ (с ИИ), ч & 7.0 & 13.73 & 34.2 & 45.6 \\
Идеализированное ускорение $S_E$ & $\times 3.43$ & $\times 4.08$ & $\times 2.22$ & $\times 0.88$ \\
Достижимое время завершения, ч & 8.4 & 16.0 & 35.2 & 52.8 \\
Достижимое ускорение $S_E^{\text{ach}}$ & $\times 2.86$ & $\times 3.50$ & $\times 2.16$ & $\times 0.76$ \\
Выполнено условие ускорения? & Да & Да & Да & Нет \\
\bottomrule
\end{tabular}
\end{table}

\begin{remark}[Final Interpretation]
\begin{enumerate}
  \item В примерах 1--3 достигается ускорение: как по идеализированной формуле, так и при проверке реалистичной раскладки неделимых задач.
  \item Различие между идеализированным и достижимым ускорением обусловлено самой длинной задачей: увеличение параллелизма не может сократить время завершения ниже её длительности.
  \item Пример 2 даёт наибольшее ускорение благодаря сочетанию $k_i < 1$ для всех задач и большего параллелизма $P=3$.
  \item Пример 3 показывает более умеренный эффект, поскольку одна из наиболее сложных задач имеет $k_i > 1$ и становится критическим ограничителем расписания.
  \item Пример 4 показывает отрицательный сценарий: даже при $P=2$ порог параллелизма не выполнен ($P < \bar{k}_w$), поэтому $T_{ai}^{(0)} > T_h$ и ускорение не достигается.
\end{enumerate}
\end{remark}

\detailedsummaryheading{Full Parameters of the Illustrative Scenario}

Перейдём от абстрактных наборов задач к иллюстративному сценарию разработки
небольшого микросервиса. Это синтетический сценарий с заданными авторами
параметрами, а не эмпирическое исследование конкретного проекта. Один и тот же проект рассматривается
в четырёх вариантах, чтобы разделить влияние двух факторов. В вариантах 1--3
при неизменном наборе задач меняется \emph{конфигурация процесса}: параллелизм и
режим работы агентов. В варианте 4 сохраняется лучшая из этих конфигураций, но
меняется \emph{декомпозиция}: крупная задача на критическом пути разбивается на
подзадачи. Такое сравнение позволяет отдельно оценить эффект организации
работы и эффект более точного разбиения проекта.

Базовая трудоёмкость единицы работы, как и ранее, принята $X = 4$~часа.

\detailedexampleheading{Scenario Definition}

Рассмотрим проект небольшого микросервиса для обработки пользовательских заявок. Сервис принимает заявку через HTTP API, проверяет входные данные, сохраняет состояние, запускает асинхронную обработку и предоставляет средства наблюдаемости. Проект достаточно компактен для полного разбора и в то же время содержит типичные элементы промышленной разработки: контракт API, хранение данных, бизнес-правила, фоновую обработку, тесты и инфраструктурную конфигурацию.

Для конкретности будем считать, что сервис должен поддерживать следующий минимальный набор возможностей:
\begin{itemize}
  \item создание заявки по HTTP-запросу;
  \item проверку обязательных полей и единый формат ошибок;
  \item сохранение заявки и её статуса в базе данных;
  \item перевод заявки по состояниям: создана, принята в обработку, обработана, завершена с ошибкой;
  \item фоновую обработку с повторными попытками при временных ошибках;
  \item структурированные журналы, метрики, проверку работоспособности и базовую конфигурацию;
  \item набор модульных и интеграционных тестов;
  \item упаковку и минимальную конфигурацию запуска в среде непрерывной интеграции и развёртывания.
\end{itemize}

Для прозрачной параметризации иллюстративного сценария используем условные единицы сложности story points (SP). Примем, что $2$~SP соответствуют одной единице $Z_i$. Тогда один story point соответствует $X/2 = 2$~часам, а стоимость единицы $Z_i$ по-прежнему равна $X = 4$~часам. Подзадачам с нечётным числом SP соответствуют дробные значения сложности ($1$~SP $= 0.5\,Z$), что согласуется с условием $Z_i \in \mathbb{Q}^{+}$ в Определении~1. В вариантах 1--3 расчёты проводятся по крупным задачам, а в варианте 4 непосредственно используется разбиение на подзадачи.

Базовое разбиение проекта на крупные задачи приведено в Таблице~\ref{tab:practical-baseline}.

\begin{table}[ht!]
\centering
\caption{Базовое разбиение проекта на задачи}
\label{tab:practical-baseline}
\begin{tabularx}{\textwidth}{@{}c>{\raggedright\arraybackslash}Xcc@{}}
\toprule
$i$ & Задача & Оценка, SP & Сложность $Z_i$ \\
\midrule
1 & API-контракт и проверка запросов & 4 & 2 \\
2 & Слой хранения (модель, миграции, репозиторий) & 6 & 3 \\
3 & Бизнес-логика (основные сценарии и переходы состояний) & 8 & 4 \\
4 & Фоновый обработчик (асинхронная обработка и повторы) & 6 & 3 \\
5 & Наблюдаемость и конфигурация (логи, метрики, health-check) & 4 & 2 \\
6 & Тесты и интеграционные проверки & 6 & 3 \\
7 & Упаковка и конфигурация развёртывания (Docker, CI) & 4 & 2 \\
\bottomrule
\end{tabularx}
\end{table}

Состав крупных задач раскрыт в Таблице~\ref{tab:practical-subtasks}. Такое уточнение важно для планирования: на уровне подзадач видно, какие работы можно независимо передавать агентам, а какие требуют заранее зафиксированного контракта или результата другой задачи. Детализация используется далее для выявления зависимостей, возникших из-за укрупнения графа, и в варианте~4, где задача тестирования разбивается на отдельные части.

\begin{table}[ht!]
\centering
\caption{Подзадачи и story points}
\label{tab:practical-subtasks}
\begin{tabularx}{\textwidth}{@{}c>{\raggedright\arraybackslash}p{0.20\textwidth}>{\raggedright\arraybackslash}Xc@{}}
\toprule
$i$ & Задача & Подзадача & SP \\
\midrule
1 & API & Описать эндпоинты, DTO и формат ошибок & 2 \\
1 & API & Реализовать проверку данных и обработку ошибок & 2 \\
2 & Хранение & Спроектировать таблицы и миграции & 2 \\
2 & Хранение & Реализовать репозиторий и методы доступа & 3 \\
2 & Хранение & Добавить базовые тестовые данные для локальной проверки & 1 \\
3 & Бизнес-логика & Описать состояния и допустимые переходы & 2 \\
3 & Бизнес-логика & Реализовать основные сценарии использования & 4 \\
3 & Бизнес-логика & Обработать ошибки доменного уровня & 2 \\
4 & Фоновый обработчик & Реализовать выборку задач в работу & 2 \\
4 & Фоновый обработчик & Добавить повторные попытки и обработку отказов & 3 \\
4 & Фоновый обработчик & Синхронизировать статусы с основной моделью & 1 \\
5 & Наблюдаемость & Добавить структурные логи и метрики & 2 \\
5 & Наблюдаемость & Настроить конфигурацию и проверку работоспособности & 2 \\
6 & Тесты & Покрыть бизнес-логику модульными тестами & 2 \\
6 & Тесты & Добавить интеграционные тесты API и хранения & 3 \\
6 & Тесты & Проверить фоновую обработку и повторные попытки & 1 \\
7 & Упаковка & Подготовить Docker/service-конфигурацию & 2 \\
7 & Упаковка & Добавить команды непрерывной интеграции и дымовой тест & 2 \\
\bottomrule
\end{tabularx}
\end{table}

Суммарная сложность проекта по крупным задачам:
\[
  \sum_{i=1}^{7} Z_i = 2 + 3 + 4 + 3 + 2 + 3 + 2 = 19.
\]
Эквивалентная сумма в story points равна $38$~SP, что при стоимости $2$~часа за SP также даёт $76$~часов базовой работы.

Время выполнения без ИИ (последовательная работа одного разработчика):
\[
  T_h = 19 \cdot 4 = 76~\text{ч}.
\]

Задачи связаны отношениями предшествования. Для планирования параллельной работы зададим граф $G=(V,E)$, вершины которого соответствуют крупным задачам из Таблицы~\ref{tab:practical-baseline}. Перечень зависимостей приведён в Таблице~\ref{tab:practical-dependencies}.

\begin{table}[ht!]
\centering
\caption{Граф зависимостей между крупными задачами}
\label{tab:practical-dependencies}
\begin{tabularx}{\textwidth}{@{}>{\raggedright\arraybackslash}p{0.29\textwidth}>{\raggedright\arraybackslash}p{0.12\textwidth}>{\raggedright\arraybackslash}X@{}}
\toprule
Задача & Зависит от & Что можно делать параллельно \\
\midrule
1. API-контракт и проверка запросов & --- & Стартовая задача; задаёт интерфейсы для остальных работ \\
2. Слой хранения & 1 & Можно выполнять параллельно с наблюдаемостью после фиксации API \\
3. Бизнес-логика & 1, 2 & Частично проектируется после API, реализация требует хранения \\
4. Фоновый обработчик & 2, 3 & Требует модели данных и правил переходов состояний \\
5. Наблюдаемость и конфигурация & 1 & Можно вести параллельно с хранением и бизнес-логикой \\
6. Тесты и интеграционные проверки & 1, 2, 3, 4 & Модульные тесты можно начинать раньше, интеграционные зависят от ядра \\
7. Упаковка и развёртывание & 1, 5 & Можно готовить параллельно после фиксации интерфейсов и конфигурации \\
\bottomrule
\end{tabularx}
\end{table}

В виде рёбер граф можно записать так:
\[
  E = \{1{\to}2,\; 1{\to}3,\; 1{\to}5,\; 1{\to}7,\; 2{\to}3,\; 2{\to}4,\; 3{\to}4,\; 3{\to}6,\; 4{\to}6,\; 5{\to}7\}.
\]

Из этого графа следует естественное волновое планирование:
\begin{enumerate}
  \item сначала фиксируется API-контракт и базовые DTO;
  \item затем параллельно можно брать хранение и наблюдаемость;
  \item после хранения становится доступна полноценная бизнес-логика, а упаковку можно готовить параллельно;
  \item фоновый обработчик зависит от бизнес-логики и модели данных;
  \item интеграционные тесты и проверки фоновой обработки завершают критический путь.
\end{enumerate}

Таким образом, в проекте сочетаются параллелизуемые участки и жёсткие зависимости. В вариантах 1--3 используются агрегированные значения $Z_i$ и коэффициенты $k_i$, однако время завершения определяется уже не балансировкой независимых задач, а расписанием на графе из Таблицы~\ref{tab:practical-dependencies}.

\paragraph{Actual Dependencies and Coarsening Artifacts.}
Укрупнённый граф скрывает важное различие. Одни рёбра Таблицы~\ref{tab:practical-dependencies} отражают \emph{содержательные} зависимости: бизнес-логика не может опираться на ещё не созданную модель данных, а фоновый обработчик --- на неопределённые переходы состояний. Другие рёбра возникают только из-за \emph{укрупнения задач}: связь проведена между задачами целиком, хотя фактически зависимы лишь отдельные подзадачи. Например, зависимость задачи 6 от задач 1--4 означает, что набор тестов содержит проверки каждого компонента. При этом модульные тесты бизнес-логики можно запускать сразу после задачи 3, интеграционные тесты --- после завершения API и слоя хранения, а только проверки повторных попыток требуют готового фонового обработчика. Сходным образом зависимость $3 \to 4$ на уровне подзадач относится лишь к модели состояний и переходов объёмом $2$~SP из $8$; остальные сценарии бизнес-логики обработчику не нужны. Содержательную зависимость нельзя устранить перепланированием, тогда как зависимость, вызванную укрупнением, можно уточнить, заменив одну вершину несколькими подзадачами. Поэтому доминирование критического пути в последующих расчётах отчасти определяется выбранной гранулярностью графа. Вариант~4 показывает этот эффект количественно.

В вариантах 1--3 проект и набор крупных задач неизменны, но различаются конфигурации процесса $\sigma = (P, r)$: число параллельных потоков $P$, режимы работы агентов $r$ и распределение задач. Поскольку коэффициент $k_i = k_i^{(r)}$ зависит от тройки «задача $\times$ инструмент $\times$ режим», интерактивная работа в варианте 1, смешанный режим в варианте 2 и делегированная работа по подробным спецификациям в варианте 3 задают разные векторы $\{k_i\}$ для одного набора задач. Инструмент и правила приёмки в сценарии фиксированы. Следовательно, сравниваются конфигурации целиком, а не изолированное изменение одного параметра. В варианте~4 воспроизводится конфигурация варианта~3 и меняется только декомпозиция, что позволяет отдельно оценить её влияние.

На уровне M4 явно учитывается конкуренция потоков за внимание разработчика. Для каждой задачи задана человеческая составляющая $h_i$, включающая постановку, проверку и направление исправлений. Расписание считается допустимым только при отсутствии пересечений между человеческими интервалами разных задач и при непрерывном закреплении одного агента за каждой начатой, но ещё не принятой задачей. Если человеческая фаза не может начаться немедленно, агент остаётся заблокированным в очереди; другие агенты продолжают работу. В интерактивном режиме $h_i = k_i$, поскольку разработчик сохраняет контекст задачи в течение всего времени её выполнения; в делегированном режиме $h_i$ заметно меньше $k_i$, но не равно нулю из-за постановки и приёмки. Значения $h_i$, как и $k_i$, заданы экспертно. Сохраняются два упрощающих предположения: $C(P) \equiv 1$, то есть межагентные интеграционные издержки не моделируются, и $\gamma(P) \equiv 1$. Влияние затрат на переключение контекста рассматривается отдельно при сравнении вариантов.

\detailedexampleheading{Variant 1. Interactive Work in a Single Stream}

Разработчик последовательно решает задачи с помощью ИИ-ассистента. Для каждой задачи он формулирует запрос, получает результат, проверяет и при необходимости исправляет его. Параллелизм отсутствует ($P = 1$). Конфигурация имеет вид $\sigma_1 = (P{=}1,\ \text{интерактивный режим})$.

\paragraph{Parameters.}
\begin{table}[ht!]
\centering
\caption{Параметры варианта 1}
\label{tab:practical-v1-params}
\begin{tabular}{clccc}
\toprule
$i$ & Задача & $Z_i$ & $k_i$ & $h_i$ \\
\midrule
1 & API-контракт и проверка запросов & 2 & 0.80 & 0.80 \\
2 & Слой хранения & 3 & 0.90 & 0.90 \\
3 & Бизнес-логика & 4 & 0.85 & 0.85 \\
4 & Фоновый обработчик & 3 & 0.95 & 0.95 \\
5 & Наблюдаемость и конфигурация & 2 & 0.75 & 0.75 \\
6 & Тесты & 3 & 0.80 & 0.80 \\
7 & Упаковка и развёртывание & 2 & 0.90 & 0.90 \\
\bottomrule
\end{tabular}
\end{table}

Коэффициенты $k_i$ отражают умеренное ускорение: ИИ помогает генерировать шаблонный код и тесты, но каждый результат проходит через ручную проверку. Для задач с высокой архитектурной сложностью (бизнес-логика, фоновый обработчик) коэффициент ближе к единице. Человеческая доля $h_i = k_i$: в интерактивном режиме разработчик удерживает контекст задачи всё её время (агентная составляющая $a_i = 0$).

\paragraph{Calculations.}
\begin{align*}
  \sum_{i=1}^{7} Z_i k_i &= 2{\cdot}0.80 + 3{\cdot}0.90 + 4{\cdot}0.85 + 3{\cdot}0.95 + 2{\cdot}0.75 + 3{\cdot}0.80 + 2{\cdot}0.90 \\
  &= 1.60 + 2.70 + 3.40 + 2.85 + 1.50 + 2.40 + 1.80 = 16.25.
\end{align*}

Средневзвешенная нормированная длительность:
\[
  \bar{k}_w = \frac{16.25}{19} \approx 0.855.
\]

Веса задач на графе равны $w_i = Z_i k_i X$:
\[
  \{w_i\}_{i=1}^{7} = \{6.4,\; 10.8,\; 13.6,\; 11.4,\; 6.0,\; 9.6,\; 7.2\}~\text{ч}.
\]

Полная работа:
\[
  W = \sum_{i=1}^{7} w_i = 16.25 \cdot 4 = 65.0~\text{ч}.
\]

Идеализированное время при независимых задачах, то есть нижняя оценка по суммарной работе при $P = 1$:
\[
  T_{ai}^{(0)} = \frac{W}{P} = \frac{65.0}{1} = 65.0~\text{ч}.
\]

Идеализированное ускорение без учёта зависимостей:
\[
  S_E = \frac{76}{65.0} \approx 1.17.
\]

Критический путь графа проходит через задачи $1 \to 2 \to 3 \to 4 \to 6$:
\[
  L = 6.4 + 10.8 + 13.6 + 11.4 + 9.6 = 51.8~\text{ч}.
\]

Нижняя оценка времени завершения с учётом графа зависимостей:
\[
  B_4 = \max\left\{\frac{W}{P},\; L,\; H\right\} = \max\{65.0,\; 51.8,\;65.0\} = 65.0~\text{ч}.
\]

При $P=1$ расписание последовательно. Задачи можно выполнить, например, в топологическом порядке $1,2,3,4,5,6,7$. Время завершения равно полному объёму работы:
\[
  T(\pi_1)=T^{*}=B_4=65.0~\text{ч}.
\]

Достижимое ускорение с учётом графа:
\[
  S_E^{\text{ach}} = \frac{76}{65.0} \approx 1.17.
\]

Человеческий ресурс (уровень M4): поскольку $h_i = k_i$,
\[
  H = X \sum_{i=1}^{7} Z_i h_i = 65.0~\text{ч} = W,
  \qquad
  \bar{h}_w = \bar{k}_w \approx 0.855,
  \qquad
  \frac{1}{\bar{h}_w} \approx 1.17.
\]

\paragraph{Interpretation.}
Интерактивный режим обеспечивает непосредственный контроль результата. ИИ снижает трудоёмкость каждой задачи ($\bar{k}_w \approx 0.855$), но отсутствие параллелизма ограничивает общий эффект. Критический путь не становится активным ограничением: при $P=1$ величина $W/P=65.0$~ч превышает $L=51.8$~ч. Ускорение составляет $\times 1.17$, что соответствует экономии примерно 11 часов из исходных 76.

На уровне M4 достигнутое ускорение $\times 1.17$ в точности совпадает с потолком по человеческому ресурсу $1/\bar{h}_w$. В интерактивном режиме время завершения равно суммарному человеческому времени $H$. Следовательно, дальнейшая реорганизация процесса не даст выигрыша без снижения $h$, то есть без передачи части работы агентам для автономного выполнения.

\detailedexampleheading{Variant 2. Interactive and Delegated Tasks}

В этом варианте доступны два агентных потока ($P = 2$). Задачи, требующие архитектурных решений, выполняются интерактивно, а более формализованные задачи делегируются агенту и принимаются разработчиком после автономной фазы. Режим относится к задаче, а не закрепляется за конкретным агентом: свободный агент может получить следующую готовую задачу любого предусмотренного режима.

Задачи распределяются по режимам следующим образом:
\begin{itemize}
  \item \textbf{Интерактивный режим}: API-контракт, бизнес-логика, фоновый обработчик --- задачи, требующие контекста и архитектурных решений.
  \item \textbf{Делегированный режим}: слой хранения, наблюдаемость, тесты, упаковка --- задачи с более формализованными требованиями.
\end{itemize}

Конфигурация: $\sigma_2 = (P{=}2,\ \rho)$ со смешанным назначением режимов $\rho$ --- задачи потока 1 выполняются в интерактивном режиме, задачи потока 2 --- в делегированном.

\paragraph{Parameters.}
\begin{table}[ht!]
\centering
\caption{Параметры варианта 2}
\label{tab:practical-v2-params}
\begin{tabular}{clcccc}
\toprule
$i$ & Задача & $Z_i$ & $k_i$ & $h_i$ & Режим \\
\midrule
1 & API-контракт и проверка запросов & 2 & 0.75 & 0.75 & интер. \\
2 & Слой хранения & 3 & 0.75 & 0.20 & делег. \\
3 & Бизнес-логика & 4 & 0.85 & 0.85 & интер. \\
4 & Фоновый обработчик & 3 & 0.90 & 0.90 & интер. \\
5 & Наблюдаемость и конфигурация & 2 & 0.70 & 0.20 & делег. \\
6 & Тесты & 3 & 0.70 & 0.20 & делег. \\
7 & Упаковка и развёртывание & 2 & 0.80 & 0.20 & делег. \\
\bottomrule
\end{tabular}
\end{table}

В делегированном режиме коэффициенты $k_i$ несколько ниже, чем в варианте~1. Это следствие смены режима: агент работает самостоятельно и не ожидает промежуточных ответов разработчика. Такой эффект проявляется уже в одном потоке и согласуется с определением $k_i^{(r)}$ как характеристики тройки «задача $\times$ инструмент $\times$ режим»; инструмент в сценарии фиксирован. Для интерактивных задач коэффициенты близки к варианту~1, а $h_i = k_i$, поскольку разработчик занят в течение всего времени выполнения. Для делегированных задач принимается $h_i = 0.20$; эта доля включает постановку и проверку результата. В расписании человеческое время делится поровну между входной спецификацией и итоговой проверкой, а агентный поток удерживается на всём интервале от начала спецификации до завершения проверки.

\paragraph{Calculations.}
\begin{align*}
  \sum_{i=1}^{7} Z_i k_i &= 2{\cdot}0.75 + 3{\cdot}0.75 + 4{\cdot}0.85 + 3{\cdot}0.90 + 2{\cdot}0.70 + 3{\cdot}0.70 + 2{\cdot}0.80 \\
  &= 1.50 + 2.25 + 3.40 + 2.70 + 1.40 + 2.10 + 1.60 = 14.95.
\end{align*}

Средневзвешенная нормированная длительность:
\[
  \bar{k}_w = \frac{14.95}{19} \approx 0.787.
\]

Веса задач на графе равны $w_i = Z_i k_i X$:
\[
  \{w_i\}_{i=1}^{7} = \{6.0,\; 9.0,\; 13.6,\; 10.8,\; 5.6,\; 8.4,\; 6.4\}~\text{ч}.
\]

Полная работа:
\[
  W = \sum_{i=1}^{7} w_i = 14.95 \cdot 4 = 59.8~\text{ч}.
\]

Идеализированное время при независимых задачах, то есть нижняя оценка по суммарной работе при $P = 2$:
\[
  T_{ai}^{(0)} = \frac{W}{P} = \frac{59.8}{2} = 29.9~\text{ч}.
\]

Идеализированное ускорение без учёта зависимостей:
\[
  S_E = \frac{76}{29.9} \approx 2.54.
\]

Критический путь графа проходит через задачи $1 \to 2 \to 3 \to 4 \to 6$:
\[
  L = 6.0 + 9.0 + 13.6 + 10.8 + 8.4 = 47.8~\text{ч}.
\]

Человеческое время (уровень M4):
\[
  \sum_{i=1}^{7} Z_i h_i = 2{\cdot}0.75 + 3{\cdot}0.20 + 4{\cdot}0.85 + 3{\cdot}0.90 + 2{\cdot}0.20 + 3{\cdot}0.20 + 2{\cdot}0.20 = 9.6,
\]
\[
  H = 9.6 \cdot 4 = 38.4~\text{ч},
  \qquad
  \bar{h}_w = \frac{9.6}{19} \approx 0.505,
  \qquad
  \frac{1}{\bar{h}_w} \approx 1.98.
\]

Нижняя оценка времени завершения с учётом графа зависимостей и человеческого ресурса:
\[
  B_4 = \max\left\{\frac{W}{P},\; L,\; H\right\} = \max\{29.9,\; 47.8,\; 38.4\} = 47.8~\text{ч}.
\]

Раскладка по потокам из Таблицы~\ref{tab:practical-v2-params}, соблюдающая зависимости:
\begin{center}
\small
\begin{tabularx}{\textwidth}{cXX}
\toprule
Поток & Интервалы задач, ч & Простой, ч \\
\midrule
1 & 1: $[0,6.0]$; 3: $[15.0,28.6]$; 4: $[28.6,39.4]$ & $[6.0,15.0]$, ожидание слоя хранения \\
2 & 2: $[6.0,15.0]$; 5: $[15.0,20.6]$; 7: $[20.6,27.0]$; 6: $[39.4,47.8]$ & $[27.0,39.4]$, ожидание зависимостей \\
\bottomrule
\end{tabularx}
\end{center}

Если учитывать только два агентных потока и отношения предшествования, эта
раскладка имеет срок $47.8$~ч и совпадает с нижней оценкой. Однако допустимым
расписанием уровня M4 она не является, как показывает следующая проверка.
\[
  T_{\mathrm{agents}}=B_4=47.8~\text{ч}.
\]

\paragraph{Human-Resource Check (Level M4).}
Приведённое расписание допустимо по потокам и зависимостям, но \emph{невыполнимо по человеческому ресурсу}. Разработчик непрерывно занят интерактивными задачами на интервале $[15.0,\,39.4]$: сначала задачей 3, затем задачей 4. В то же время делегированные задачи 5 $[15.0,\,20.6]$ и 7 $[20.6,\,27.0]$ требуют по $1.6$~ч участия человека для постановки и проверки. Эти интервалы пересекаются с интерактивной работой, что невозможно при одном разработчике. Следовательно, формально допустимое по агентным потокам расписание не предусматривает ресурса для своевременной приёмки делегированных результатов.

Конфликт устраним перестройкой расписания: делегированные задачи и их человеческие фазы размещаются в окнах, когда разработчик свободен. При этом начатая задача непрерывно удерживает назначенного агента, включая ожидание итоговой проверки. Скорректированное расписание:
\begin{itemize}
  \item Агент 1: задача 1 $[0,\,6.0]$; задача 2 $[6.0,\,15.0]$; задача 3 $[15.0,\,28.6]$; задача 4 $[28.6,\,39.4]$; задача 6 $[39.4,\,47.8]$.
  \item Агент 2: задача 5 $[7.2,\,12.8]$; задача 7 $[12.8,\,41.4]$.
  В задаче 7 спецификация занимает $[12.8,\,13.6]$, автономная работа ---
  $[13.6,\,18.4]$, ожидание проверки --- $[18.4,\,40.6]$, проверка ---
  $[40.6,\,41.4]$.
  \item Человек: задача 1 $[0,\,6.0]$; спецификация задачи 2 $[6.0,\,7.2]$; спецификация задачи 5 $[7.2,\,8.0]$; проверка задачи 5 $[12.0,\,12.8]$; спецификация задачи 7 $[12.8,\,13.6]$; проверка задачи 2 $[13.8,\,15.0]$; задача 3 $[15.0,\,28.6]$; задача 4 $[28.6,\,39.4]$; спецификация задачи 6 $[39.4,\,40.6]$; проверка задачи 7 $[40.6,\,41.4]$; проверка задачи 6 $[46.6,\,47.8]$. Интервалы человека попарно не пересекаются, суммарная занятость равна $38.4$~ч $= H$.
\end{itemize}
Все зависимости соблюдены. Задача 3 начинается в момент $15.0$ после полной приёмки задачи 2; задача 7 --- после приёмки задачи 5; задача 6 --- после задачи 4. Период ожидания задачи 7 не используется для другой работы на агенте 2. Несмотря на эту локальную блокировку, время завершения сохраняется равным $47.8$~ч, поскольку задача 7 находится вне критического пути. Таким образом, человеческий ресурс не увеличивает срок только при явном планировании фаз и очереди; исходное расписание было невыполнимым.

Скорректированное расписание допустимо и достигает нижней оценки, поэтому оно
оптимально:
\[
  T(\pi_2)=T^{*}=B_4=47.8~\text{ч}.
\]

Достижимое ускорение с учётом графа и человеческого ресурса:
\[
  S_E^{\text{ach}} = \frac{76}{47.8} \approx 1.59.
\]

\paragraph{Interpretation.}
Если считать задачи независимыми, нижняя оценка $W/P=29.9$~ч соответствовала бы ускорению примерно $\times 2.54$. Однако граф зависимостей существенно сокращает доступный параллелизм: задачи $1 \to 2 \to 3 \to 4 \to 6$ образуют критический путь длиной $47.8$~ч. С учётом графа достижимое ускорение снижается до $\times 1.59$. Наблюдаемость и развёртывание действительно выполняются параллельно, но находятся вне критического пути и почти не влияют на общий срок.

Ключевое наблюдение: параллелизм $P = 2$ помогает только там, где граф допускает независимую работу. В этом варианте большая часть трудоёмкости оказывается свернута в последовательную цепочку, поэтому балансировка потоков без учёта зависимостей завышает эффект конфигурации в иллюстративном сценарии.

Человеческий ресурс пока не является активным ограничением: $H = 38.4$~ч $< L = 47.8$~ч, а достигнутое ускорение ниже потолка $1/\bar{h}_w \approx 1.98$ ($1.59 < 1.98$). Однако этот резерв реализуется только при явном планировании окон проверки. Три интерактивные задачи с $h = k$ занимают разработчика в общей сложности на $30.4$~ч, поэтому проверки делегированных задач необходимо размещать в оставшихся промежутках.

\detailedexampleheading{Variant 3. Delegated Agents with Detailed Specifications}

В этом варианте четыре ИИ-агента работают параллельно ($P = 4$). Конфигурация задаётся как $\sigma_3 = (P{=}4,\ \text{делегированный режим с подробными спецификациями})$. Коэффициенты $k_i$ ниже, чем в предыдущих вариантах, благодаря смене профиля режима, а не увеличению параллелизма: агенты получают подробные спецификации и автономно выполняют агентные фазы, после чего результат принимает разработчик. Граф зависимостей между крупными задачами остаётся прежним и ограничивает доступный параллелизм.

\paragraph{Parameters.}
\begin{table}[ht!]
\centering
\caption{Параметры варианта 3}
\label{tab:practical-v3-params}
\begin{tabular}{clccc}
\toprule
$i$ & Задача & $Z_i$ & $k_i$ & $h_i$ \\
\midrule
1 & API-контракт и проверка запросов & 2 & 0.65 & 0.25 \\
2 & Слой хранения & 3 & 0.70 & 0.25 \\
3 & Бизнес-логика & 4 & 0.80 & 0.25 \\
4 & Фоновый обработчик & 3 & 0.75 & 0.25 \\
5 & Наблюдаемость и конфигурация & 2 & 0.60 & 0.25 \\
6 & Тесты & 3 & 0.65 & 0.25 \\
7 & Упаковка и развёртывание & 2 & 0.70 & 0.25 \\
\bottomrule
\end{tabular}
\end{table}

Для всех задач принята одинаковая человеческая составляющая $h_i = 0.25$. Участие разработчика состоит в подготовке подробной спецификации и последующей проверке результата; реализация полностью передаётся агенту.

\paragraph{Calculations.}
\begin{align*}
  \sum_{i=1}^{7} Z_i k_i &= 2{\cdot}0.65 + 3{\cdot}0.70 + 4{\cdot}0.80 + 3{\cdot}0.75 + 2{\cdot}0.60 + 3{\cdot}0.65 + 2{\cdot}0.70 \\
  &= 1.30 + 2.10 + 3.20 + 2.25 + 1.20 + 1.95 + 1.40 = 13.40.
\end{align*}

Средневзвешенная нормированная длительность:
\[
  \bar{k}_w = \frac{13.40}{19} \approx 0.705.
\]

Веса задач на графе равны $w_i = Z_i k_i X$:
\[
  \{w_i\}_{i=1}^{7} = \{5.2,\; 8.4,\; 12.8,\; 9.0,\; 4.8,\; 7.8,\; 5.6\}~\text{ч}.
\]

Полная работа:
\[
  W = \sum_{i=1}^{7} w_i = 13.40 \cdot 4 = 53.6~\text{ч}.
\]

Идеализированное время при независимых задачах, то есть нижняя оценка по суммарной работе при $P = 4$:
\[
  T_{ai}^{(0)} = \frac{W}{P} = \frac{53.6}{4} = 13.4~\text{ч}.
\]

Идеализированное ускорение без учёта зависимостей:
\[
  S_E = \frac{76}{13.4} \approx 5.67.
\]

Критический путь графа проходит через задачи $1 \to 2 \to 3 \to 4 \to 6$:
\[
  L = 5.2 + 8.4 + 12.8 + 9.0 + 7.8 = 43.2~\text{ч}.
\]

Человеческое время (уровень M4):
\[
  \sum_{i=1}^{7} Z_i h_i = 0.25 \cdot 19 = 4.75,
  \qquad
  H = 4.75 \cdot 4 = 19.0~\text{ч},
  \qquad
  \bar{h}_w = 0.25,
  \qquad
  \frac{1}{\bar{h}_w} = 4.0.
\]

Нижняя оценка времени завершения с учётом графа зависимостей и человеческого ресурса:
\[
  B_4 = \max\left\{\frac{W}{P},\; L,\; H\right\} = \max\{13.4,\; 43.2,\; 19.0\} = 43.2~\text{ч}.
\]

Одна допустимая раскладка по четырём агентам, соблюдающая зависимости:
\begin{itemize}
  \item Агент 1: задача 1 $[0,5.2]$, задача 2 $[5.2,13.6]$, задача 3 $[13.6,26.4]$, задача 4 $[26.4,35.4]$, задача 6 $[35.4,43.2]$.
  \item Агент 2: задача 5 $[6.7,11.5]$, задача 7 $[15.6,21.2]$.
  \item Агенты 3 и 4 простаивают, потому что при таком крупном графе критический путь доминирует над доступным параллелизмом.
\end{itemize}

Это расписание достигает нижней оценки:
\[
  T(\pi_3)=T^{*}=B_4=43.2~\text{ч}.
\]

\paragraph{Human-Resource Check (Level M4).}
Чтобы агент не начинал работу до постановки задачи, человеческое время делится поровну между подготовкой входной спецификации и итоговой проверкой; между этими этапами агент выполняет реализацию. Блоки разработчика размещаются так: задача 1 --- $[0,\,1.0]$ и $[4.2,\,5.2]$; задача 2 --- $[5.2,\,6.7]$ и $[12.1,\,13.6]$; задача 5 --- $[6.7,\,7.7]$ и $[10.5,\,11.5]$; задача 3 --- $[13.6,\,15.6]$ и $[24.4,\,26.4]$; задача 7 --- $[15.6,\,16.6]$ и $[20.2,\,21.2]$; задача 4 --- $[26.4,\,27.9]$ и $[33.9,\,35.4]$; задача 6 --- $[35.4,\,36.9]$ и $[41.7,\,43.2]$. Интервалы попарно не пересекаются, а каждая входная спецификация предшествует агентной части соответствующей задачи. Суммарная занятость человека равна $19.0$~ч $= H$ при длительности проекта $43.2$~ч, поэтому расписание выполнимо без перестройки.

Достижимое ускорение с учётом графа и человеческого ресурса:
\[
  S_E^{\text{ach}} = \frac{76}{43.2} \approx 1.76.
\]

\paragraph{Interpretation.}
При четырёх агентах суммарная работа снижается до $W=53.6$~ч, а оценка для независимых задач составляет $W/P=13.4$~ч. Однако при укрупнённом графе это значение недостижимо: цепочка $1 \to 2 \to 3 \to 4 \to 6$ имеет длину $43.2$~ч и становится главным ограничением. Поэтому после учёта зависимостей ускорение составляет не $\times 5.67$, а $\times 1.76$.

Улучшение относительно варианта~2 следует интерпретировать осторожно. Критический путь сокращается с $47.8$ до $43.2$~ч за счёт уменьшения коэффициентов $k_i$ при переходе к делегированному режиму с подробными спецификациями. Сам рост $P$ с 2 до 4 не влияет на $L$: дополнительные агенты не ускоряют последовательную цепочку крупных задач и в данном расписании простаивают. Следующее изменение сценария состоит не в дальнейшем увеличении числа агентов, а в разбиении крупных задач критического пути на меньшие, независимо проверяемые части. Именно такое изменение рассматривается в варианте~4.

Уровень M4 задаёт дополнительную верхнюю границу. В данном варианте потолок по человеческому ресурсу $1/\bar{h}_w = 4.0$ заметно выше достигнутого ускорения $\times 1.76$, поэтому активным ограничением остаётся $L$. Тем не менее при выбранном режиме ($\bar{h}_w = 0.25$) никакая декомпозиция и никакое число агентов не обеспечат ускорение выше $\times 4.0$: 19 часов последовательного человеческого времени устранить нельзя. По мере сокращения критического пути активным ограничением станет $H$.

\detailedexampleheading{Variant 4. The Same Configuration with Finer Decomposition}

От варианта~3 этот вариант отличается только декомпозицией. Число агентов $P = 4$, делегированный режим с подробными спецификациями, коэффициенты $k$ и человеческие доли $h$ сохраняются. Меняется лишь граф: крупная задача 6 «Тесты» заменяется тремя подзадачами из Таблицы~\ref{tab:practical-subtasks}, для каждой из которых задаются точные зависимости. В вариантах 1--3 конфигурация менялась при фиксированном разбиении; здесь, напротив, разбиение меняется при фиксированной конфигурации. Поэтому различие с вариантом~3 обусловлено только декомпозицией.

\paragraph{Parameters.}
Задачи 1--5 и 7 не меняются (Таблица~\ref{tab:practical-v3-params}). Задача 6 дробится согласно Таблице~\ref{tab:practical-subtasks}:

\begin{table}[ht!]
\centering
\caption{Подзадачи задачи 6 в варианте 4}
\label{tab:practical-v4-params}
\begin{tabular}{clcccl}
\toprule
$i$ & Подзадача & SP & $Z_i$ & $k_i$ / $h_i$ & Зависит от \\
\midrule
6а & Модульные тесты бизнес-логики & 2 & 1 & 0.65 / 0.25 & 3 \\
6б & Интеграционные тесты API и хранения & 3 & 1.5 & 0.65 / 0.25 & 1, 2 \\
6в & Проверка фоновой обработки и повторов & 1 & 0.5 & 0.65 / 0.25 & 4 \\
\bottomrule
\end{tabular}
\end{table}

Предполагается, что разбиение \emph{не меняет объём работы}: подзадачи наследуют $k = 0.65$ и $h = 0.25$ исходной задачи. Поэтому величины $W$, $\bar{k}_w$ и $H$ сохраняются, а сравнение с вариантом~3 выделяет структурный эффект. На практике декомпозиция требует дополнительных затрат на спецификацию, постановку и приёмку каждой подзадачи. По принятому правилу учёта они входят в $k^{(r)}$ и при чрезмерном дроблении увеличивают его. С другой стороны, для небольших задач с чёткими границами $k$ может снижаться. Таким образом, неизменность $k$ и $h$ --- нейтральное сценарное предположение, а чувствительность результата к издержкам декомпозиции требует отдельной калибровки.

Сложности подзадач дробные ($Z_{6\text{б}} = 1.5$, $Z_{6\text{в}} = 0.5$): нечётное число SP при курсе $2$~SP $= 1\,Z$. Веса подзадач:
\[
  w_{6\text{а}} = 1 \cdot 0.65 \cdot 4 = 2.6~\text{ч}, \qquad
  w_{6\text{б}} = 1.5 \cdot 0.65 \cdot 4 = 3.9~\text{ч}, \qquad
  w_{6\text{в}} = 0.5 \cdot 0.65 \cdot 4 = 1.3~\text{ч};
\]
в сумме $7.8$~ч --- вес исходной задачи 6.

\paragraph{Calculations.}
Агрегаты совпадают с вариантом~3:
\[
  W = 53.6~\text{ч}, \qquad
  \bar{k}_w \approx 0.705, \qquad
  \frac{W}{P} = 13.4~\text{ч}, \qquad
  H = 19.0~\text{ч}, \qquad
  \frac{1}{\bar{h}_w} = 4.0.
\]

Меняется критический путь. Монолитная вершина 6 и входящие в неё рёбра $3 \to 6$, $4 \to 6$ заменены тремя вершинами с точными зависимостями $3 \to 6\text{а}$, $\{1,2\} \to 6\text{б}$, $4 \to 6\text{в}$, и самым длинным путём становится
\[
  1 \to 2 \to 3 \to 4 \to 6\text{в}: \qquad
  L = 5.2 + 8.4 + 12.8 + 9.0 + 1.3 = 36.7~\text{ч}
\]
(конкурирующие пути короче: $1 \to 2 \to 3 \to 6\text{а}$ даёт $29.0$~ч, $1 \to 2 \to 6\text{б}$ --- $17.5$~ч, $1 \to 5 \to 7$ --- $15.6$~ч). Хвост критического пути сократился с $7.8$~ч (вся задача 6) до $1.3$~ч (только та её часть, которая действительно ждёт фонового обработчика).

Нижняя оценка времени завершения:
\[
  B_4 = \max\left\{\frac{W}{P},\; L,\; H\right\} = \max\{13.4,\; 36.7,\; 19.0\} = 36.7~\text{ч}.
\]

Допустимая раскладка, соблюдающая зависимости:
\begin{itemize}
  \item Агент 1: задача 1 $[0,\,5.2]$, задача 2 $[5.2,\,13.6]$, задача 3 $[13.6,\,26.4]$, задача 4 $[26.4,\,35.4]$, задача 6в $[35.4,\,36.7]$.
  \item Агент 2: задача 5 $[6.7,\,11.5]$; для задачи 7 входная спецификация занимает $[19.5,\,20.5]$, агентная реализация --- $[20.5,\,24.1]$, после чего результат ожидает освобождения разработчика до проверки $[28.4,\,29.4]$. На интервале ожидания $[24.1,\,28.4]$ агент остаётся закреплён за задачей 7 и не может получить другую работу.
  \item Агент 3: задача 6б $[15.6,\,19.5]$, задача 6а $[27.9,\,30.5]$. Агент 4 простаивает.
\end{itemize}

Это расписание достигает нижней оценки:
\[
  T(\pi_4)=T^{*}=B_4=36.7~\text{ч}.
\]

\paragraph{Human-Resource Check (Level M4).}
Как и в варианте~3, человеческое время каждой задачи делится поровну между входной спецификацией и итоговой проверкой. Блоки разработчика: задача 1 --- $[0,\,1.0]$ и $[4.2,\,5.2]$; задача 2 --- $[5.2,\,6.7]$ и $[12.1,\,13.6]$; задача 5 --- $[6.7,\,7.7]$ и $[10.5,\,11.5]$; задача 3 --- $[13.6,\,15.6]$ и $[24.4,\,26.4]$; задача 6б --- $[15.6,\,16.35]$ и $[18.75,\,19.5]$; задача 7 --- $[19.5,\,20.5]$ и $[28.4,\,29.4]$; задача 4 --- $[26.4,\,27.9]$ и $[33.9,\,35.4]$; задача 6а --- $[27.9,\,28.4]$ и $[30.0,\,30.5]$; задача 6в --- $[35.4,\,35.65]$ и $[36.45,\,36.7]$. Интервалы попарно не пересекаются, каждая входная спецификация предшествует агентной реализации, а суммарная занятость человека равна $19.0$~ч $= H$. Следовательно, расписание допустимо по человеческому ресурсу и действительно достигает $36.7$~ч.

Достижимое ускорение с учётом графа и человеческого ресурса:
\[
  S_E^{\text{ach}} = \frac{76}{36.7} \approx 2.07.
\]

\paragraph{Interpretation.}
Время завершения сократилось с $43.2$ до $36.7$~ч, а ускорение выросло с $\times 1.76$ до $\times 2.07$, то есть на $18\%$, без добавления агентов, смены режима или улучшения инструмента. Эффект получен за счёт замены одной крупной вершины тремя подзадачами, уже выделенными в Таблице~\ref{tab:practical-subtasks}. Зависимость «тестирование после завершения всего ядра» оказалась следствием укрупнения графа; её уточнение сократило критический путь на $6.5$~ч. В рассматриваемом сравнении результат определяется декомпозицией, а не числом агентов.

Только переход от варианта 3 к варианту 4 представляет собой сравнение при прочих равных: сокращение времени с $43.2$ до $36.7$~ч обусловлено декомпозицией при неизменных $P$, $k$ и $h$. При переходе от варианта 2 к варианту 3 одновременно меняются $P$ и режим, поэтому их влияние не разделено. Тем не менее расчёт показывает, что уменьшение активной границы $L$ с $47.8$ до $43.2$~ч связано со снижением $k$ на критическом пути. Само значение $P$ в формулу $L$ не входит, а дополнительные агенты простаивают.

Сравнительный вывод для вариантов 3 и 4 сохраняется и без точной абсолютной
калибровки. Согласно Предложению~\ref{prop:invariance}, умножение всех агентных
и человеческих фаз обоих вариантов на один общий множитель $\kappa>0$ изменит
оба времени завершения в $\kappa$ раз, но оставит неизменными отношение
$43.2/36.7$ и ранжирование вариантов. Эта устойчивость не распространяется на
неодинаковые ошибки по задачам или на изменение относительных профилей $k_i$ и
$h_i$.

Декомпозицию можно продолжить итеративно. Следующий кандидат --- зависимость $3 \to 4$, которая на уровне подзадач сводится к модели состояний объёмом $2$~SP из $8$. Если выделить её в отдельную раннюю подзадачу, фоновый обработчик можно начинать до завершения остальных сценариев бизнес-логики. Предел такого дробления уже заметен: путь $1 \to 2 \to 3 \to 4$ длиной $35.4$~ч почти полностью определяет $L = 36.7$~ч, а снизу время ограничено человеческим ресурсом $H = 19.0$~ч и потолком ускорения $\times 4.0$. По мере дальнейшего сокращения критического пути активным ограничением станет $H$, как предсказывает уровень M4.

\detailedexampleheading{Comparison of Variants}

\begin{table}[ht!]
\centering
\caption{Сравнение вариантов организации работы с учётом графа зависимостей}
\label{tab:practical-comparison}
\small
\begin{tabularx}{\textwidth}{Xcccc}
\toprule
Показатель & Вариант 1 & Вариант 2 & Вариант 3 & Вариант 4 \\
\midrule
Параллелизм $P$ & 1 & 2 & 4 & 4 \\
Режим & интерактивный & смешанный & делег.\ со спец. & делег.\ со спец. \\
Декомпозиция & крупная & крупная & крупная & дроблёные тесты \\
$W$, ч & 65.0 & 59.8 & 53.6 & 53.6 \\
$W/P$, ч & 65.0 & 29.9 & 13.4 & 13.4 \\
$L$, ч & 51.8 & 47.8 & 43.2 & 36.7 \\
$H$, ч & 65.0 & 38.4 & 19.0 & 19.0 \\
$B_4$, ч & 65.0 & 47.8 & 43.2 & 36.7 \\
$T(\pi)$, ч & 65.0 & 47.8 & 43.2 & 36.7 \\
$S_E^{\text{ach}}$ & $\times 1.17$ & $\times 1.59$ & $\times 1.76$ & $\times 2.07$ \\
Потолок $1/\bar{h}_w$ & $\times 1.17$ & $\times 1.98$ & $\times 4.00$ & $\times 4.00$ \\
\bottomrule
\end{tabularx}
\end{table}

Из таблицы видно, что после учёта зависимостей итоговое ускорение определяется не только суммарной работой $W$ и числом агентов $P$, но и длиной критического пути $L$ --- а сама длина $L$ зависит от выбранной декомпозиции.

\begin{enumerate}
  \item \textbf{Нижняя оценка $W/P$ полезна, но оптимистична.} Для варианта~2 она равна $29.9$ ч, а для варианта~3 --- $13.4$ ч. Эти числа достижимы только при достаточно независимых задачах и согласованных по времени запросах к человеку; ожидание в очереди увеличивает занятость потоков сверх $W$.

  \item \textbf{В вариантах 2--4 результат определяет критический путь.} В вариантах 2 и 3 цепочка $1 \to 2 \to 3 \to 4 \to 6$ задаёт время завершения $47.8$ и $43.2$~ч соответственно. После дробления тестов в варианте 4 критический путь сокращается до $36.7$~ч, но остаётся активным ограничением. Поэтому достижимое ускорение составляет $\times 1.59$, $\times 1.76$ и $\times 2.07$, а не значения, рассчитанные в предположении полной независимости задач.

  \item \textbf{Увеличения $P$ самого по себе недостаточно.} При переходе от $P=2$ к $P=4$ оценка $W/P$ для независимых задач резко снижается, но общее время меняется умеренно: дополнительные агенты не могут начать задачи критического пути до завершения их предшественников.

  \item \textbf{Эффект декомпозиции выделен сравнением при прочих равных.} Переходы $1 \to 2$ и $2 \to 3$ одновременно меняют несколько параметров конфигурации. Только при переходе $3 \to 4$ сохраняются $P$, режим, $k$ и $h$, поэтому рост ускорения с $\times 1.76$ до $\times 2.07$ относится к изменению декомпозиции и не требует дополнительных ресурсов.

  \item \textbf{Человеческий ресурс задаёт потолок и ограничивает расписание.} При уменьшении человеческой составляющей потолок $1/\bar{h}_w$ возрастает: $\times 1.17 \to \times 1.98 \to \times 4.00$. В варианте 1 он уже достигнут, и процесс ограничен человеком. В вариантах 2--4 остаётся резерв, но использовать его можно только при явном планировании окон обслуживания. В перестроенном варианте 2 агент задачи 7 заблокирован в очереди $22.2$~ч, однако это не увеличивает makespan, поскольку задача находится вне критического пути.

  \item \textbf{Умеренные затраты на переключение не меняют ранжирование примера.} При $\gamma(P)=1+0.1(P-1)$ точный пересчёт с уточнённой временной сеткой даёт (часы):
  \[
  \begin{array}{c|ccccc}
  \text{вариант} & W_4 & L_4 & \gamma H & B_4 & T^* \\ \hline
  2 & 63.640 & 51.320 & 42.240 & 51.320 & 51.36 \\
  3 & 59.300 & 47.700 & 24.700 & 47.700 & 47.70 \\
  4 & 59.300 & 40.450 & 24.700 & 40.450 & 40.45
  \end{array}
  \]
  Поэтому сохраняется порядок $4<3<2$, а human-ветвь остаётся неактивной. При более консервативной оценке $h_i=0.5k_i$ в вариантах 3--4 получаем $H=26.8$~ч и $\gamma(4)H\approx34.8$~ч; для варианта 4 это уже близко к критическому пути. Следовательно, по мере сокращения $L_4$ человеческий ресурс может стать активным ограничением.
\end{enumerate}

\paragraph{Implications of the Illustrative Scenario.}
\begin{enumerate}
  \item \textbf{Зависимости следует учитывать явно.} Балансировка независимых задач завышает ускорение, если работы связаны графом зависимостей. Для рассмотренного примера на уровне M4 нижняя оценка равна $\max\{W/P, L, H\}$ при $C(P) = \gamma(P) = 1$.

  \item \textbf{Критический путь стал главным ограничением.} В вариантах~2--4 время завершения совпадает с $L$. При этом состав критического пути зависит от гранулярности графа: после дробления задачи тестирования в его хвост входит только проверка повторных попыток, тогда как модульные и интеграционные тесты выполняются раньше и параллельно.

  \item \textbf{Рост числа агентов не даёт пропорционального ускорения.} Если граф содержит крупные последовательно связанные задачи, увеличение $P$ быстро перестаёт влиять на результат: часть агентов простаивает в ожидании обязательных предшественников.

  \item \textbf{Декомпозиция даёт измеримый эффект.} Разбиение одной задачи тестирования сократило время завершения на $6.5$~ч и повысило ускорение с $\times 1.76$ до $\times 2.07$ без изменения $P$, $k$ и $h$. Следующие кандидаты следует искать среди рёбер критического пути, которые связывают крупные задачи целиком, хотя фактическая зависимость относится лишь к одной подзадаче.

  \item \textbf{Планировать необходимо не только агентов, но и участие человека.} Проверка делегированных задач выполняется одним разработчиком последовательно, поэтому соответствующие интервалы и очередь к ним должны быть явно включены в расписание уровня M4. Ожидающий агент сохраняет свой поток; передать ему другую задачу до приёмки текущей нельзя. Длительная непрерывная интерактивная работа блокирует обслуживание параллельных результатов, а предельное ускорение ограничивается величиной $1/\bar{h}_w$ независимо от числа агентов.
\end{enumerate}

\ifdefined\FULLARTICLE
  \let\finishdocument\relax
\else
  \def\finishdocument{\end{document}}
\fi
\finishdocument
